\subsection{Cambiamenti di base ed operatori unitari}
\[
\ket{\psi} = \sum_i \ket{e_i} \braket{e_i|\psi} = \sum_i \ket{e_i'}\braket{e_i'|\psi}
\]
\[
\braket{e_i'|\psi} = \sum_j \braket{e_i'|e_j} \braket{e_j|\psi}
\]
e chiamo
\[
U_{ij} = \braket{e_i'|e_j}
\]
e \(U\) \`e la matrice di cambiamento di base, con
\[
U = \sum_m \ket{e_m} \bra{e_m'}
\]

Inoltre, se \(U \ket{\psi} = \ket{\psi'}\), vale che
\begin{enumerate}
\item \(\braket{e_i|\psi'} = \sum_m \braket{e_i|e_m'} \braket{e_m'|\psi} = \braket{e_i'|\psi}\)
\item \(U_{ij} = \sum_m \braket{e_i|e_m} \braket{e_m'|e_j} = \sum_m \braket{e_i'|e_m'} \braket{e_m'|e_j'} = \braket{e_i'|U|e_j'} = \braket{e_i'|e_j}\)
\end{enumerate}

\(U\) \`e unitario se
\[
U^\dag = U^{-1}, \qquad U^\dag U = \mathbbm{1}
\]
e scrivo
\[
U = \sum_m \ket{e_m} \bra{e_m'}
\]
\[
U^\dag = \sum_m \ket{e_m'} \bra{e_m}
\]
\[
U U^\dag = \sum_{n,m} \ket{e_n} \braket{e_n'|e_m'} \bra{e_m} = \sum_m \ket{e_m} \bra{e_m} = \mathbbm{1} = U^\dag U
\]
Gli operatori unitari sono matrici di cambio base.

Prendo due stati \(\ket{\psi}\) e \(\ket{\varphi}\), allora per \(U\) unitario
\[
\ket{\psi'} = U \ket{\psi} \qquad \ket{\varphi'} = U \ket{\varphi}
\]
\[
\braket{\psi'|\varphi'} = \braket{\psi|UU^\dag|\varphi} = \braket{\psi|\varphi}
\]
cio\`e i cambiamenti di base attraverso operatori unitari mantiene invariato il prodotto scalare (si conserva la probabilit\`a). Ho a che fare con la reversibilit\`a: in particolare, l'operatore di misura non \`e reversibile.
\[
\ket{\psi} \xrightarrow{\text{evoluzione}} U \ket{\psi} \qquad U \text{ operatore unitario}
\]
\[
\ket{\psi} \xrightarrow{\text{misura}} \ket{e_k} = N \, P_k \ket{\psi} \qquad P_k \text{ operatore di proiezione}
\]
Infatti gli operatori di proiezione non sono unitari perch\'e se esistesse l'inverso \(P^{-1}\) avrei \(P^2 = P\), cio\`e \(P = \mathbbm{1}\).